% Options for packages loaded elsewhere
\PassOptionsToPackage{unicode}{hyperref}
\PassOptionsToPackage{hyphens}{url}
\documentclass[
  icelandic,
]{book}
\usepackage{xcolor}
\usepackage{amsmath,amssymb}
\setcounter{secnumdepth}{5}
\usepackage{iftex}
\ifPDFTeX
  \usepackage[T1]{fontenc}
  \usepackage[utf8]{inputenc}
  \usepackage{textcomp} % provide euro and other symbols
\else % if luatex or xetex
  \usepackage{unicode-math} % this also loads fontspec
  \defaultfontfeatures{Scale=MatchLowercase}
  \defaultfontfeatures[\rmfamily]{Ligatures=TeX,Scale=1}
\fi
\usepackage{lmodern}
\ifPDFTeX\else
  % xetex/luatex font selection
\fi
% Use upquote if available, for straight quotes in verbatim environments
\IfFileExists{upquote.sty}{\usepackage{upquote}}{}
\IfFileExists{microtype.sty}{% use microtype if available
  \usepackage[]{microtype}
  \UseMicrotypeSet[protrusion]{basicmath} % disable protrusion for tt fonts
}{}
\makeatletter
\@ifundefined{KOMAClassName}{% if non-KOMA class
  \IfFileExists{parskip.sty}{%
    \usepackage{parskip}
  }{% else
    \setlength{\parindent}{0pt}
    \setlength{\parskip}{6pt plus 2pt minus 1pt}}
}{% if KOMA class
  \KOMAoptions{parskip=half}}
\makeatother
\usepackage{longtable,booktabs,array}
\usepackage{calc} % for calculating minipage widths
% Correct order of tables after \paragraph or \subparagraph
\usepackage{etoolbox}
\makeatletter
\patchcmd\longtable{\par}{\if@noskipsec\mbox{}\fi\par}{}{}
\makeatother
% Allow footnotes in longtable head/foot
\IfFileExists{footnotehyper.sty}{\usepackage{footnotehyper}}{\usepackage{footnote}}
\makesavenoteenv{longtable}
\usepackage{graphicx}
\makeatletter
\newsavebox\pandoc@box
\newcommand*\pandocbounded[1]{% scales image to fit in text height/width
  \sbox\pandoc@box{#1}%
  \Gscale@div\@tempa{\textheight}{\dimexpr\ht\pandoc@box+\dp\pandoc@box\relax}%
  \Gscale@div\@tempb{\linewidth}{\wd\pandoc@box}%
  \ifdim\@tempb\p@<\@tempa\p@\let\@tempa\@tempb\fi% select the smaller of both
  \ifdim\@tempa\p@<\p@\scalebox{\@tempa}{\usebox\pandoc@box}%
  \else\usebox{\pandoc@box}%
  \fi%
}
% Set default figure placement to htbp
\def\fps@figure{htbp}
\makeatother
\ifLuaTeX
\usepackage[bidi=basic]{babel}
\else
\usepackage[bidi=default]{babel}
\fi
% get rid of language-specific shorthands (see #6817):
\let\LanguageShortHands\languageshorthands
\def\languageshorthands#1{}
\setlength{\emergencystretch}{3em} % prevent overfull lines
\providecommand{\tightlist}{%
  \setlength{\itemsep}{0pt}\setlength{\parskip}{0pt}}
\usepackage[]{biblatex}
\addbibresource{book.bib}
\usepackage{booktabs}
\ExecuteBibliographyOptions{sorting=nyt,sortlocale=is-IS}
\linespread{1.2}
\usepackage{bookmark}
\IfFileExists{xurl.sty}{\usepackage{xurl}}{} % add URL line breaks if available
\urlstyle{same}
\hypersetup{
  pdftitle={Nótur um stærðfræðimenntun},
  pdfauthor={Ingólfur Gíslason},
  pdflang={is-IS},
  hidelinks,
  pdfcreator={LaTeX via pandoc}}

\title{Nótur um stærðfræðimenntun}
\author{Ingólfur Gíslason}
\date{2026-01-12}

\begin{document}
\maketitle

{
\setcounter{tocdepth}{1}
\tableofcontents
}
\chapter{Um textann}\label{umtextann}

Þessi texti er verk í vinnslu. Hann er um stærðfræðimenntun, einkum kennslu og nám (frekar en kerfis- eða félagsfræðilega sýn). Hún er hugsuð sem kennsluefni í námskeiðinu Stærðfræði í grunnskóla sem er fyrir stærðfræðikennaranema við Menntavísindasvið Háskóla Íslands. Hún gæti þó nýst nemendum í öðrum námskeiðum eða starfandi kennurum. Textinn gerir ráð fyrir að lesendur hafi bók Henri Picciotto og Robin Pemantle, \emph{There Is No One Way to Teach Math} \autocite{picciotto2025} til hliðsjónar og hann á að þjóna sem stuðningur og dýpkun við lestur og notkun þeirrar bókar.

Í textanum er áherslan á undirstöðuhugtök og praktískar leiðir í kennslu, auk rýni í stærðfræðilegt efni grunnskólans. Gengið er út frá því að lesendur vinni bæði saman og með kennara, en áhugasamar manneskjur utan slíks samhengis gætu ef til vill nýtt sér textann líka.

\section{Notkun bókar og tungumál}\label{notkun-buxf3kar-og-tungumuxe1l}

Ég ávarpa lesanda textans ýmist sem nemanda (hann), lesanda (hann), en stundum nota ég „þau``, „öll``, og önnur orð til að gefa til kynna að öll eru velkomin að lesa og læra af þessum texta, og persónur í verkefnatextum geta verið af ólíku kyni.

Bókina má aðlaga og nýta að vild, að hluta eða í heild, með eftirfarandi skilyrðum:

\begin{enumerate}
\def\labelenumi{\arabic{enumi}.}
\item
  Bókin sé ekki nýtt til að valda manneskjum skaða, græða peninga, eða til að stuðla að auknum ójöfnuði milli fólks.
\item
  Ef umtalsverðir hlutar eru nýttir í öðru verki sé upprunans hér getið.
\end{enumerate}

\chapter{Skilningur}\label{skilningur}

Hvað er skilningur? Hvað er skilningur í stærðfræði?

\section{Venslaskilningur og tækniskilningur}\label{venslaskilningur-og-tuxe6kniskilningur}

Við getum rifjað upp tvígreiningu Richard Skemps í venslaskilning og tækniskilning \autocite{skemp1976}. Í örstuttu máli drögum við þessi hugtök saman:

\begin{itemize}
\tightlist
\item
  Venslaskilningur er að vita hvernig hægt er að setja nota stærðfræðihugtök og -aðferðir og vita hvers vegna þær eru viðeigandi og hvers vegna þær skila réttum niðurstöðum.
\item
  Tækniskilningur er að vita hvernig á að reikna tiltekin dæmi án þess að vita hvers vegna þau eru reiknuð á þann hátt, hvort það sé vit í reikningunum, eða hvernig aðferðin tengist öðrum aðferðum eða hugtökum.
\end{itemize}

Í bók Picciotto og Pemantle er \emph{skilningur} nokkurn veginn það sama og Skemp kallar venslaskilning.

\section{Skilningur í daglegu máli og í stærðfræði}\label{skilningur-uxed-daglegu-muxe1li-og-uxed-stuxe6ruxf0fruxe6uxf0i}

Til frekari dýptar getum við velt fyrir okkur hvað orðið skilningur eiginlega merkir. Í skáldsögu Fríðu Ísberg, \emph{Merking} \autocite{isberg2021} (annað mikilvægt hugtak!) segir:

\begin{quote}
Íslenska orðið „skilningur`` hefur ákveðna einangrunarmerkingu; við reynum að skilja eitt frá öðru - hið ljóta frá hinu fallega, manninn frá dýrinu - til að einangra það og eyða svo því sem við viljum ekki. En latneska orðið „comprehendere`` - sem enskan notar - er andstæða þess: það þýðir að taka saman. „Com-`` er svipað forskeyti og „sam-`` á íslensku og „prehendere`` þýðir að grípa, ná tökum á einhverju. (bls. 194)
\end{quote}

Fleiri rannsakendur hafa lagt áherslu á að skilningur felist bæði í að aðskilja, greina sundur, og líka að taka saman, ná utan um, grípa um. Mjög fróðleg umfjöllun um skilning í íslensku og í sögulegu ljósi kennslubóka í stærðfræði er í grein Kristínar Bjarnadóttur, \emph{Er áhersla á skilning nýmæli?} \autocite{bjarnadottir2022}

Til að skýra betur hvað það merkir að \emph{skilja hugtak} segja Picciotto og Pemantle að það að skilja hugtak ætti yfirleitt að þýða að nemandi geti:

\begin{itemize}
\tightlist
\item
  Útskýrt það, útskýra \emph{hvers vegna} eitthvað er eins og það er, ekki bara nefna orð. Við ættum að biðja nemendur reglulega að útskýra bæði munnlega og skriflega.
\item
  Að snúa ferlum við. Þú skilur ekki dreifireglu nema þú getir þáttað líka, þú skilur ekki jöfnur nema þú getir búið til (margar) jöfnur sem hafa lausnina \(4\), og getir búið til jöfnu út frá grafi. Það að geta snúið við ferlum (reversibility) er prófsteinn á skilning, leið til að dýpka skilning, og í sumum tilvikum önnur leið að skilningi.
\item
  Svegjanleg notkun margra leiða. Að skilja jöfnur: með prófun, með gröfum, með töflum, með tækni, auk bókstafareiknings.
\item
  Að geta \emph{tengt milli ólíkra framsetninga}, til dæmis tákna (algebrustæða), töflu (gildistöflu) og grafs (teikning í hnitakerfi).
\item
  Yfirfærsla á ný samhengi. Til dæmis að tengja hlutföll við einslaga myndir, og fjarlægð milli punkta við reglu Pýþagórasar.
\item
  Að vita hvenær það á ekki við. Dæmi: ekki eru öll sambönd línuleg. Þess vegna er mikilvægt að rannsaka gagndæmi, hvenær eitthvað á ekki við.
\end{itemize}

\subsection{Verkefni}\label{verkefni}

Hugsum okkur spurninguna

\begin{quote}
Af hverju er \(2(x+3) = 2x + 6\)?
\end{quote}

Hvaða ályktanir um skilning nemanda getum dregið út frá eftirfarandi „svörum``, og hvernig gætum við brugðist við svörunum til að kanna skilninginn betur:

\begin{enumerate}
\def\labelenumi{\arabic{enumi}.}
\tightlist
\item
  „Það er út af dreifireglunni?{}``
\item
  „Til dæmis ef \(x=1\) þá þetta bara \(2(1+3)=2 \cdot 4\) og það er alveg eins hægt að skipta \(4\) niður og margfalda hvern hlut fyrir sig, \(2 \cdot 1 = 2\) og \(2 \cdot 3 = 6\) og leggja saman, \(2+6=8\).
\item
  „Ef ég skoða rétthyrning með eina hlið \(2\) og hina \(x+3\) og reikna flatarmálið, þá skiptir ekki máli hvort ég reikna hvern bút fyrir sig og legg saman eða allan í einu`` (með fylgir eftirfarandi teikning):
\end{enumerate}

\begin{figure}
\centering
\pandocbounded{\includegraphics[keepaspectratio]{_main_files/figure-latex/reitir-1.pdf}}
\caption{\label{fig:reitir}Rétthyrningur með reitum af breidd \(x\) og \(3\) og hæð \(2\).}
\end{figure}

\chapter{Talnaskyn, táknskyn, aðgerðaskyn}\label{skyn}

Orðið skyn á hér að fela í sér bæði skilning og skynjun.

\section{Talnaskyn (number sense)}\label{talnaskyn-number-sense}

Með talnaskyni er meðal annars átt við að hve miklu leyti við getum reiknað í huganum, tengt saman talnastaðreyndir sem við þekkjum og áttað okkur á stæðarsamanburði talna. Ef við höfum gott talnaskyn, höfum við „tengslahugsun`` um tölur og reikning, og notum okkur staðreyndir um tölur og talnasambönd til að leysa verkefni, og við áttum okkur á samanburði talna (líka brota) út frá eiginleikum (en þurfum ekki endilega að „reikna þær út``, sjáum til dæmis að \(\frac{9}{14}\) er örugglega stærri tala en \(\frac{1}{2}\), vegna þess að \(9>7\)).

Í rannsóknum á stærðfræðikennslu byggða á skilningi barna (cognitively guided instruction) kemur í ljós (Sjá \autocite{kristinsdottir2021}):

\begin{quote}
að börn sem fá tækifæri til að leysa þrautir með eigin aðferðum þróa með sér lausnaleiðir sem ákvarðast af fyrri reynslu þeirra og þroskastigi. Þetta ferli er líkt hjá öllum börnum og skiptist í þrjú meginstig:
\end{quote}

\begin{quote}
Stig 1 --- hlutbundið líkan sem fylgir söguþræði
Stig 2 --- talning
Stig 3 --- tengslahugsun og nýting talnastaðreynda
\end{quote}

Hlutbundið líkan geta verið hlutir, fingur eða skýringarmyndir. Á stigi 2 er líkanið (eins og á stigi 1) í huganum -- börnin þurfa ekki lengur að hafa hlutbundna framsetningu á tölu. Það þýðir þó ekki að þau noti ekki til dæmis fingur eða skýringamyndir. Á stigi 3 getum við unnið með tölur á sveigjanlegan hátt og nýtt staðreyndir um tölur og talnasambönd.

\subsection{Dæmi}\label{duxe6mi}

\begin{itemize}
\tightlist
\item
  Ef við reiknum \(13-9\) með því að hugsa \(13 = 10 + 3\) og notum okkur að \(10-9=1\) til að sjá að \(13-9 = 10 - 9 + 3 = 1 + 3 = 4\).
\item
  Ef við reiknum \(4 \cdot 13 = 2 \cdot 2 \cdot 13 = 2 \cdot 26 = 52\)
\end{itemize}

\section{Táknskyn (symbol sense)}\label{tuxe1knskyn-symbol-sense}

Með táknskyni er meðal annars átt við:

\begin{enumerate}
\def\labelenumi{\alph{enumi}.}
\tightlist
\item
  Að geta valið hvort gagnlegt sé að nota tákn (eins og bókstafi fyrir breytur og óþekktar tölur) við lausn verkefnis eða ekki
\item
  Að geta túlkað tákn sem notuð eru til að lýsa reikningum eða aðstæðum
\item
  Fimi í bókstafareikningi (að umbreyta táknarunum í aðrar jafngildar táknarunur)
\item
  Val á réttum táknum fyrir tilteknar aðstæður
\end{enumerate}

Nemendur þurfa að átta sig á því að bókstafir fyrir breytur gegna ólíkum hlutverkum en hafa þó tiltekna eiginleika. Til dæmis: Breytur \ldots{}

\begin{itemize}
\tightlist
\item
  geta táknað tölur sem vantar (og hægt er að finna út hver er)
\item
  fylgja öllum reiknireglum sem gilda um tölur
\item
  eru stundum tengdar öðrum breytum
\item
  hafa stundum eitt tiltekið gildi, en stundum geta þær haft hvaða gildi sem er
\end{itemize}

\subsection{Verkefni}\label{verkefni-1}

Eftirfarandi spurningar reyna á táknskyn með ólíkum hætti. Í hverri línu er fullyrðing sem segja á hvort sé stundum, alltaf eða aldrei sönn. Við höfum áhuga á skilningi svo hér þarf að útskýra \emph{hvers vegna} í sérhverjum lið. Við viljum líka benda á undir hvaða flokk táknskyns hér að ofan hvert dæmi reynir helst á.

\begin{itemize}
\tightlist
\item
  \(x+2 < 2x\).
\item
  \(-x\) er neikvæð tala.
\item
  Það skiptir máli hvort við skrifum \(4a+10\) eða \(4b+10\).
\item
  Ef \(a\) og \(b\) eru breytur, þá er ómögulegt að \(a=b\).
\item
  \((x+5)^2 = x^2 + 5^2\).
\item
  \(-(y-1) = -y -1\) ?
\item
  \(-a^2 = (-a)^2\) ?
\item
  Ef þú margfaldar þrjár tölur sem standa saman í talnaröðinni (eins og \(13\), \(14\) og \(15\)) verður útkoman alltaf margfeldi af \(6\) (með öðrum orðum: \(6\) gengur upp í henni).
\item
  Ef \(3\) gengur upp í þversummu tölu þá gengur \(3\) gengur upp í tölunni.
\end{itemize}

Táknskyn er ekki eitthvað sem kemur fljótt eða sjálfkrafa. Það þarf að kenna nemendum að vinna með tákn.

\section{Aðgerðaskyn (operation sense)}\label{auxf0geruxf0askyn-operation-sense}

Sameiginlegt með talnaskyni og táknskyni er aðgerðaskyn. Þá er átt við að skilja reikniaðgerðir, ólíka eiginleika þeirra og tengsl þeirra á milli, auk þess að hafa vit á því að nota þær við viðeigandi aðstæður og verkefnum.

\subsection{Dæmi}\label{duxe6mi-1}

Skoðum rununa \(5, 8, 11, 14, 17...\) Hvað er í gangi hér?

\begin{itemize}
\tightlist
\item
  Endurtekin samlagning, \(5, 5+3, 5+3+3, ...\) sem hægt er að tjá með margföldun
\item
  Táknað með \(5 + 3n\) og tengist framsetningu á línu \(y = 3x+5\)
\end{itemize}

Við viljum meðal annars að nemendur skilji að endurtekna samlagningu má reikna með margföldun, og átti sig á og geti notað dreifiregluna \(a(b+c)\) í sinni talnahugsun og hugarreikningi.

\subsection{Verkefni}\label{verkefni-2}

Víxlregla um samlagningu segir að jafnan \(a+b = b+a\) sé alltaf sönn.

\begin{enumerate}
\def\labelenumi{\arabic{enumi}.}
\tightlist
\item
  Gildir víxlregla um frádrátt? Er hægt að segja eitthvað um tölurnar \(a\) og \(b\) ef \(a-b = b-a\)?
\item
  Er einhver regla á því hvað gerist ef liðum er víxlað í frádrætti? Hvað gerist?
\item
  Gildir víxlregla um margföldun?
\item
  Gildir víxlregla um deilingu?
\end{enumerate}

\subsection{Verkefni: Börn leysa verkefni um margföldun}\label{verkefni-buxf6rn-leysa-verkefni-um-margfuxf6ldun}

Skoðið greinina \href{https://skolathraedir.is/2021/04/21/eg-leysi-stundum-vandamalid-med-svona-hringjum-hugsun-barna-um-margfoldun/}{„Ég leysi stundum vandamálið með svona hringjum`` - Hugsun barna um margföldun}. Beinum athygli okkar að talnaskyni barnanna.

\begin{enumerate}
\def\labelenumi{\alph{enumi}.}
\item
  Horfið á Óðin leysa verkefnið „Sólveig keypti 4 poka af gulrótum. Í hverjum poka voru 6 gulrætur. Hvað keypti Sólveig margar gulrætur samtals?{}`` \href{https://youtu.be/rqlU-PLoxOs?si=U6lG57jxB6EKkLsn}{Youtube-hlekkur}. Gerið grein fyrir talnaskyni og aðgerðaskyni í lausn Óðins. Til dæmis: Hvaða talnasambönd þekkir hann og hvaða eiginleika margföldunar og samlagningar nýtir hann sér?
\item
  Horfið á Torfa leysa sama verkefni \href{https://youtu.be/OEZPAQegt9o?si=7VloBBZ5PemM2v_u}{Youtube-hlekkur}. Hvaða hæfni nýtir hann sér og hvað bendir til þess að talnaskyn hans sé á þessum punkti þróaðra heldur en hjá Óðni? (Hafið vel í huga að við erum ekki að bera saman nemendurna sjálfa heldur einungis hvað kemur fram í þessum tilteknu lausnum.)
\end{enumerate}

\subsection{Verkefni}\label{verkefni-3}

Fyllið inn í eftirfarandi töflu. Það er búið að gera fyrstu röðina. Skýringamyndir eiga að sýna inntak verkefnisins. Útreikningur á að sýna eina margföldun eða deilingu.

\begin{longtable}[]{@{}
  >{\raggedright\arraybackslash}p{(\linewidth - 6\tabcolsep) * \real{0.2500}}
  >{\raggedright\arraybackslash}p{(\linewidth - 6\tabcolsep) * \real{0.2500}}
  >{\raggedright\arraybackslash}p{(\linewidth - 6\tabcolsep) * \real{0.2500}}
  >{\raggedright\arraybackslash}p{(\linewidth - 6\tabcolsep) * \real{0.2500}}@{}}
\toprule\noalign{}
\begin{minipage}[b]{\linewidth}\raggedright
\textbf{Dæmi}
\end{minipage} & \begin{minipage}[b]{\linewidth}\raggedright
\textbf{Útskýrðu hvernig á að leysa dæmið og teiknaðu mynd til að hjálpa}
\end{minipage} & \begin{minipage}[b]{\linewidth}\raggedright
\textbf{Útreikningur}
\end{minipage} & \begin{minipage}[b]{\linewidth}\raggedright
\textbf{Svar}
\end{minipage} \\
\midrule\noalign{}
\endhead
\bottomrule\noalign{}
\endlastfoot
\textbf{1.} Ég kaupi fjórar samstæður (fernur) af þremur jógúrtdollum. Hve margar dollur af jógúrt kaupi ég alls? & Ég reikna þrjá hópa af fjórum. & 3 × 4eða4 × 3 & \textbf{12} \\
\textbf{2.} Tveimur pítsum er jafnt skipt á milli fimm manns. Hversu mikið fær hver einstaklingur? & & & \\
\textbf{3.} Max sker köku í þrjá jafna hluta. Hann borðar hálfan hluta. Hve stóran hluta (brot) af kökunni borðar hann? & & & \\
\textbf{4.} Búðu til þetta dæmi sjálfur! & & 3 ÷ 1/2 & \\
\textbf{5.} Búðu til þetta dæmi sjálfur! & & 1/2 ÷ 1/4 & \\
\end{longtable}

Hér er listi af algengum vandræðum. Verkefnið er að útskýra í hverju vandræðin felast, og stinga upp á spurningum eða vísbendingum fyrir nemanda sem lendir í vandanum.

\begin{enumerate}
\def\labelenumi{\arabic{enumi}.}
\tightlist
\item
  Nemandinn reiknar \(5 \div 2\) í spurningu 2.
\item
  Nemandinn gefur svarið \(1 \frac{1}{2}\) í spurningu 4.
\item
  Nemandinn gefur svarið \(\frac{1}{8}\) í spurningu 4.
\item
  Nemandinn reiknar \(\frac{1}{3} \div \frac{1}{2}\) í spurningu 3.
\item
  Nemandinn skrifar \(1 \div 3 \div 2\) í spurningu 3.
\item
  Nemandinn á erfitt með að gera skýringamynd.
\item
  Nemandinn klárar verkefnið og þarf útvíkkun og áskorun.
\end{enumerate}

\chapter{Kennslunálganir og stærðfræðistríð}\label{kennslunuxe1lganir-og-stuxe6ruxf0fruxe6uxf0istruxeduxf0}

Fólk er oft mjög ósammála um það hvernig stærðfræðikennsla á að fara fram.

\section{Togstreitur}\label{togstreitur}

Það sem stundum er kallað \emph{hefðbundin stærðfræðikennsla} einkennist af því að kennarinn sýnir nemendum hvernig hægt er að leysa tilteknar gerðir af verkefnum (eða „hvernig á`` að leysa þær) og nemendur æfa sig svo á samsvarandi dæmum hver fyrir sig. Í slíku fyrirkomulagi er hætt við því að áherslan verði of mikil á aðferðir án skilnings (þ.e. nemendur öðlast eingöngu tækniskilning). Markmiðið með stærðfræðikennslu er hins vegar skilningur (þ.e. venslaskilningur). Góðir stærðfræðikennarar sem vinna innan þessa ramma, eða nýta sér hann ásamt öðrum aðferðum, útskýra hlutina vel og þar er möguleiki á því að nemendur bæti við venslaskilning sinn. Góðar útskýringar geta hins vegar blekkt nemendur til að halda að þeir skilji eitthvað sem þeir gera í raun og veru ekki og það að reyna einfaldlega að muna útskýringarnar virkar ekki vel fyrir flesta. Það þarf meira til, einhverja vitsmunalega vinnu, að brjóta heilann, að hugsa á einbeittan hátt.

Þess vegna leggur stærðfræðimenntafræði til að kennslan felist meira í því að nemendur sjálfir reyni að leysa verkefni án þess að reyna að herma eftir því sem kennarinn hefur gert. Þá er hugmyndin að nemendur sjálfir geti uppgötvað stærðfræðileg sannindi og hugtök. Slíkt fyrirkomulag krefst þó vandaðs undirbúnings, vals á verkefnum og skipulags um verkefnavinnuna. Ef það er ekki gert er hætt við að nemendur læri hvorki aðferðir né nái skilningi. Venslaskilningur kemur ekki fram sjálfkrafa þó að nemendur fáist við verðug verkefni. Uppgötvanir þarf að festa í sessi með umræðu og ítrekun. Það þarf að leiða í ljós með nemendum hvað sé mikilvægt, hvað sé þess virði að festa í minni, hvernig uppgötvunin tengist öðru efni og hvernig það bregður nýju ljósi á fyrri þekkingu, auk þess sem nemendur þurfa æfingu í að beita því sem þau uppgötvuðu, eða því sem leitt var í ljós með samtali eftir heilabrotin.

Undirliggjandi í allri umræðu um stærðfræðikennslu er togstreita milli þess hve mikið kennarinn á að sýna nemendum beint (bein kennsla, kynningar, fyrirlestrar) og hve mikið nemendur geta byggt upp stærðfræðina sjálf. Hægt er að lesa um það með því að fletta upp „math wars`` á netinu. Ein af þeim röksemdum sem „hefðarsinnar`` beita er að segja að nemendur þurfi að hafa vald á grunnfærni áður en þeir geta byggt upp hugtakaskilning. Þetta er ekki rétt, en rétta myndin er ekki heldur að það sé hægt að hafa góðan hugtakaskilning án færni. Hvoru tveggja þarf til og oftast þarf að fara ítrekað á milli: bæði að reikna og að hugsa um merkingu og rökin á bakvið.

Hve mikið á að hjálpa nemendum með stærðfræðiverkefni? Það er freistandi fyrir kennara að einfaldlega sýna nemendum hvernig hægt er að leysa verkefnið. Það er líka algengt að nemendur sæki það mjög stíft. En þá er hætt við að nemendur hætti að hugsa og brjóta heilann, og læri því lítið af verkefninu. Listin er að svara þannig að við stöðvum ekki hugsun nemenda heldur hjálpum þeim að hugsa áfram. Oft er hægt að svara með spurningu á móti en einstaka sinnum er hægt að gefa vísbendingu án þess að gefa svarið eða stela allri hugsuninni. Góð vísbending beinir athygli nemandans að lykileiginleikum þess stærðfræðilega viðfangs, þeirra hugtaka eða þeirrar framsetningar sem hann er að reyna að átta sig á. Einnig getur verið nauðsynlegt að rifja upp með nemandanum merkingu einhverra orða eða tákna. Komið verður mun meira inn á þessi atriði síðar.

Önnur skyld togstreita er milli umhyggju fyrir nemendum og umhyggju fyrir greininni, eða „commitment to students`` á móti „commitment to the discipline`` eins og Picciotto og Pemantle kalla það. Samkvæmt þeim felst góð kennsla \textbf{alls ekki} í því að finna milliveg heldur \textbf{verður að gera hvoru tveggja}. Hvert tiltekið námsefni eða efnisatriði þarf sína sérstöku nálgun og „dans`` milli glímu (hugsunar, vitsmunalegrar vinnu) og skýringa og umræðu. Þetta fer fram og til baka.

Oft er of mikið gert úr því að stærðfræði byggi alltaf á því sem undan er komið eins og í beinni keðju þar sem engan hlekk má vanta. Vissulega \emph{er} stærðfræði meira í þessa veru en margar aðrar greinar en þetta er samt of einfölduð mynd. Það er oft hægt að skilja stærðfræðileg efni þó að eitthvað vanti upp á þekkingu á einhverju sem „ætti að koma á undan``.

\subsection{Dæmi}\label{duxe6mi-2}

Margföldunartaflan er sígilt dæmi um efnisatriði sem fólk deilir um. Annars vegar telja sumir rétt að leggja mikla áherslu á að nemendur æfi sig og læri \(10x10\)-töfluna utanbókar og geti fundið hvert margfeldi í henni hratt og örugglega. Hins vegar segja aðrir að það sé mikilvægara að þróa talna- og aðgerðaskyn sitt þannig að maður skilji eiginleika margföldunar vel, og geti nýtt sér það til að finna margfeldin örugglega en ef til vill ekki alveg sjálfvirkt. Fólk í þeim hópi myndi yfirleitt frekar vilja tala um að „festa í minni`` margföldunartöfluna, sem getur gerst með því að vinna með margföldun á marga vegu, heldur en að læra hana utanbókar. Það er vissulega mikilvægt að hafa hraðan og öruggan aðgang í huganum að margföldunarstaðreyndum innan \(10x10\)-töflunnar. Ef þú hefur hana ekki á þokkalegu valdi þínu verður erfitt að fylgjast með og halda áfram stærðfræðinámi vegna þess að til að skilja þáttun þarftu að þekkja þáttun einhverra talna (segjum talna \(<100\)). Ef þú sérð töluna \(24\) og veist ekki (nánast sjálfvirkt) að \(24 = 2 \cdot 12, 3 \cdot 8, 4 \cdot 6\) þá mun of mikil orka fara í að finna þetta út. Þetta þýðir ekki að besta leiðin til þess að festa margföldunartöfluna í minni sé að þylja hana utanbókar. Og það eru margar gagnlegar margföldunarstaðreyndir sem fólk öðlast ekki endilega vitund um þó að það kunni að þylja margföldunartöfluna, til dæmis:

\begin{itemize}
\tightlist
\item
  \(0 \cdot a = 0\)
\item
  \(1 \cdot a = a\)
\item
  \(ab = ba\)
\item
  \(a(bc)=a(bc)\)
\item
  \(2a\) endar á sléttri tölu
\item
  \(5a\) endar á \(5\) eða \(0\),
\item
  \(ab = 2a \cdot \frac{1}{2}b\),
\item
  \(a(b+c) = ab + bc\).
\end{itemize}

Þriðja togstreitan sem vert er að nefna hér er á milli þess að hreyfa kennsluna áfram með nýju efni og þess að rifja upp og skoða aftur eldra efni. Það þarf hvoru tveggja: hreyfingu áfram og nýtt efni en líka endurlit, upprifjun og að skoða aftur eldra efni í nýju ljósi (í ljósi nýrrar þekkingar).

\section{Hvað virkar ekki}\label{hvauxf0-virkar-ekki}

Það er hægt að útfæra stærðfræðikennslu með ólíkum hætti. Ef rökrænt samhengi greinarinnar verður nemendum ljóst, með meiri áherslu á hugtök en aðferðir og þeir takast á við krefjandi verkefni um alhæfingar og röksemdir þá skiptir minna máli hvernig kennsluaðferðir eru notaðar.

Hér eru nokkrar leiðir sem gefast yfirleitt ekki vel:

\begin{itemize}
\tightlist
\item
  Ofuráhersla á \emph{form} - tiltekin orð og rithátt.
\item
  Að samþykkja ekki svör sem ekki hafa verið einfölduð eða sett á „rétt form``.
\item
  Að læra utanbókar án skilnings.
\item
  Ég útskýri og svo æfið þið ykkur á dæmum.
\item
  Opnið bókina, skoðið skýringar og sýnidæmi og hermið eftir þeim.
\item
  Hér er nýtt verkefni (sem ekki hefur verið útskýrt) og nú finnið þið út úr því án samræðu eða nokkurs inngrips kennara.
\end{itemize}

Þó geta verið tilteknar aðstæður þar sem þessir hlutir koma að gagni.

\section{Að festa í sessi (institutionalization)}\label{auxf0-festa-uxed-sessi-institutionalization}

Fæstir geta bæði tekið glósur og hugsað (um stærðfræði) samtímis. Megnið af tíma kennslustundar ætti að fara í það að nemendur leysi verkefni.

Eftir að nemendur hafa glímt við verkefni er mikilvægt að taka saman lærdóminn, að festa í sessi það sem nemendur eiga að taka með sér úr tímanum og tengja við stærðfræði sem fræðigrein, tungumál stærðfræðinnar, hefðbundin tákn og orð yfir hlutina og leiðir til þess að sannreyna og staðfesta niðurstöður.

Þá getur kennarinn líka \emph{sagt nemendum} hvað þau eiga að skrifa (glósa).

Í lokasamantektinni þarf kennari að

\begin{itemize}
\tightlist
\item
  Draga fram lykilhugtökin skýrt og greinilega.
\item
  Skýra hvað sé mikilvægast og þess virði að muna og skrifa.
\item
  Að hjálpa nemendum að tengja nýja efnið við fyrra efni og annað efni.
\item
  Að sýna hefðbundna táknun og orð.
\end{itemize}

Þessi samantekt er samræða kennarans við allan bekkinn en hún byggir á vinnu nemenda.

\printbibliography

\end{document}
